\slajdid{03-s004}
\begin{frame}[label=03-s004]
\gusttekst
Ono što je bitno jeste da su se u jednom pakovanju \underline{po pravilu} \underline{nalazila kola iste vrste,} \underline{na primer samo dvoulazna NI} \underline{kola, ili samo invertori itd\ldots} Znači ako nam za sintezu kombinacione mreže treba dva dvoulazna NI kola, jedno troulazno NI, i jedan invertor morali smo kupiti tri čipa: sa dvoulaznim NI kolima, sa troulaznim NI kolima i sa invertorima.

\alert{Ono što je takođe jako bitno jeste da se nekorišćeni ulazi i ulazi u nekorišćene delove komponente moraju postaviti na neaktivne nivoe. U suprotnom može doći do neželjenih efekata po rad aktivnog dela mreže. Kao što se vidi nije svejedno kako ćemo od početnih logičkih funkcija uraditi sintezu kombinacione mreže.}

U prethodnom primerima nadam se da ste uočili da sa dvoulaznim NI kolom možemo napraviti kratkim spajanjem ulaznih priključaka invertor – ne treba nam čip sa invertorima. Isto tako sa dva dvoulazna NI kola možemo napraviti jedno troulazno, normalno po cenu većeg kašnjenja \(F=A\cdot B\cdot C=A\cdot(B\cdot C)\), Znači trebalo bi nam i manji broj čipova, a i kupovali bi čipove iste vrste. Opet, ako smemo da dozvolimo, kašnjenje će biti povećano pošto umesto jednog kašnjenja kroz troulazno logičko kola imamo kašnjenje kroz dva dvoulazna logička kola.
\end{frame}
